% Status: DRAFT -- check surrounding flow before direct reuse.
% Source: experiments/17_pseudo_formalization_related_work/paper_summary_and_contrast.md

\paragraph{Pseudo-formal proof verification vs.\ judge certification.}
\citet{barkallah2026pseudoformalization} propose \textsc{Pseudo-Formalization} plus Block Verification, which rewrites a natural-language mathematical proof into modules with explicit premises, conclusions, and local proof text before asking an LLM to verify each module.
The method targets a complementary bottleneck: long informal proofs are difficult for a verifier to inspect as one monolithic object.
Our protocol instead targets the verifier itself.
Rather than changing the proof representation, we test whether a Lean-aware theorem/specification judge obeys identity, bug-monotonicity, specification-monotonicity, and repeat-stability checks before its scores support benchmark claims.
Thus pseudo-formalization makes proofs easier to judge, whereas property certification asks whether the chosen judge is reliable enough to use.

\paragraph{Pseudo-formal proof verification vs.\ code-to-Lean autoformalization.}
\citet{barkallah2026pseudoformalization} introduce a structured natural-language proof format and a block verifier that improves error localization on competition and research-level math proofs.
This is closest to \textsc{VeriBench} in its scalable-oversight motivation, but the evaluated object is different.
PF+BV starts from a mathematical proof and keeps the intermediate representation in natural language, while \textsc{VeriBench} starts from executable Python and scores whether an agent produces Lean~4 code, specifications, and proofs that compile and cover gold obligations.
The contrast is therefore not informal versus formal prose style, but verifier target: PF+BV improves LLM verification of a proof argument, whereas \textsc{VeriBench} evaluates end-to-end source-code formalization.

\paragraph{Pseudo-formal proof verification vs.\ differential source grounding.}
\citet{barkallah2026pseudoformalization} decompose natural-language mathematical proofs into premise/conclusion/proof modules and verify each block with an LLM.
That decomposition is useful evidence that localizing a verification task can improve LLM oversight.
\textsc{VeriBench-DT} localizes a different relation: whether a generated Lean artifact is behaviorally grounded in the original source program.
Its adapters and differential tests define executable observations across Python and Lean, while theorem coverage handles semantic claims that cannot be reduced to sampled outputs.
PF+BV has no executable source-program oracle; \textsc{VeriBench-DT} has no need to rewrite an informal proof into pseudo-formal natural language.

\paragraph{Pseudo-formal proof verification vs.\ repository-scale dependency formalization.}
\citet{barkallah2026pseudoformalization} verify mathematical proof arguments by translating them into structured natural-language modules.
\textsc{VeriBench-Deps} instead evaluates repository-scale code formalization, where the hard boundary is dependency disposition: which internal callees must be formalized, which internal effects must be axiomatized, and which external dependencies are legitimately outside the source tree.
The shared idea is modular verification.
The difference is the module type: PF+BV modules are proof-argument blocks, while VeriBench-Deps modules are cross-file software functions with explicit formalize/axiomatize decisions.
